\section{Introduction}

\label{sec:introduction}
% ══════════════════════════════════════════════════════════════════════════════

Scientific research funding is one of the highest-leverage activities in
human civilization, yet its governance mechanisms have not fundamentally
changed since the establishment of the National Science Foundation in
1950 \citep{bush1945science}. The system exhibits several well-documented
failure modes:

\begin{enumerate}[label=(\roman*)]
  \item \textbf{Conservatism bias}: Grant panels favor incremental
    research over transformative ideas, as novel proposals lack the
    preliminary data that reviewers demand \citep{boudreau2016looking}.
  \item \textbf{Geographic concentration}: 80\% of global R\&D spending
    occurs in 10 countries \citep{unesco2021science}, excluding
    researchers in developing nations from funding opportunities.
  \item \textbf{Slow cycles}: Grant application-to-funding timelines of
    6--18 months are misaligned with fast-moving research areas
    \citep{alberts2014rescuing}.
  \item \textbf{Peer review opacity}: Anonymous peer review enables
    arbitrary gatekeeping and bias \citep{lee2013bias}.
  \item \textbf{Misaligned incentives}: Researchers are rewarded for
    publication quantity rather than knowledge quality, producing a
    replication crisis \citep{ioannidis2005most}.
  \item \textbf{IP monopolization}: Publicly funded research generates
    proprietary patents, limiting knowledge diffusion
    \citep{heller1998can}.
\end{enumerate}

\subsection{Decentralized Science (DeSci)}

The DeSci movement applies blockchain and decentralized governance
principles to scientific research \citep{hamburg2019desci}. Key tenets
include open access to research outputs, transparent funding allocation,
community governance of research priorities, and token-incentivized
contribution. Projects like VitaDAO (longevity research), LabDAO
(computational biology), and Molecule (IP-NFTs for drug discovery)
have demonstrated the viability of decentralized research coordination.

Zach Kelling extends DeSci principles to conservation science,
environmental research, and AI for biodiversity---domains that
particularly benefit from global coordination and community governance.

\subsection{Contributions}

\begin{enumerate}
  \item A \textbf{complete governance framework} for decentralized
    research funding, encompassing proposal submission, review, voting,
    funding, and accountability.

  \item \textbf{Quadratic voting} adapted for research grant allocation,
    with formal analysis showing superior preference aggregation.

  \item A \textbf{sybil-resistant reputation system} that credentials
    researchers based on verified contributions.

  \item \textbf{Incentivized peer review} through reviewer staking,
    where reviewers commit tokens to their assessments.

  \item \textbf{Retroactive public goods funding} that rewards
    high-impact research after outcomes are demonstrated.

  \item \textbf{Empirical evaluation} from 18 months of Zoo DAO
    operation comparing outcomes to traditional funding mechanisms.
\end{enumerate}


% ══════════════════════════════════════════════════════════════════════════════
